\section{Introduction}
\label{sec:intro}

Puno Quechua (ISO 639-3: \texttt{qxp}) belongs to the Southern Quechua branch
and is spoken by approximately 465{,}000 people in the Puno region of
Peru~\cite{huaman2026punoasr}. Formal education in the region is conducted
almost entirely in Spanish, so most speakers are fluent in Quechua but cannot
read or write it. That asymmetry is what makes speech recognition the
appropriate technology here rather than a convenience: a text interface excludes
precisely the people the technology is meant to serve, while a speech interface
meets them where their literacy already is.

Speech resources for Quechua exist for other varieties
\cite{cardenas2018siminchik,paccotacya2022collao}, but until recently none was
dedicated to this one. Huaman et
al.~\cite{huaman2026punoasr} released the first through a participatory
campaign: 66 hours of recordings, a benchmark across four architectures, and
open checkpoints. We start from their best system and take part in Task~1 of
the SIMBig~2026 shared task, which evaluates Puno Quechua recognition on
held-out audio in two domains, scripted and spontaneous speech.

The task ranks systems by a mean rank over \emph{four} numbers: word and
character error rate in each of the two domains. This is not a detail of
scoring. The baseline paper reports two configurations of the same model, and
averaging the four numbers separates them sharply. The first reaches 1.19 WER
on scripted speech but 13.6 WER on spontaneous speech: excellent in one domain,
unusable in the other. The second reaches 2.11 and 3.15 on those same two
measures, giving up a little scripted accuracy for a great deal of spontaneous
accuracy. The mean ranks the second above the first, and it is the
configuration we start from. Throughout this paper we therefore report all
four numbers and their mean, and we treat an improvement that moves one of them
in the wrong direction as no improvement at all.

This paper makes three contributions:

\begin{enumerate}
  \item A system that improves the published baseline from 13.15 to 9.14 mean
        WER measured on the same split with the same reference text
        (Section~\ref{sec:results}).
  \item The first quantification of sentence overlap in the established
        evaluation split for this variety: hearing a sentence read by other
        speakers is worth 1.1 WER points, a 21\% relative reduction
        (Section~\ref{sec:overlap}).
  \item Three reproducible negative results on ensembling, regularisation and
        language model decoding (Section~\ref{sec:negative}).
\end{enumerate}

% =====================================================================
